\chapter{Conclusion}
\label{chap:conclusion}
\setlength{\parskip}{0.8em}

\section{Summary of Research}
This thesis described the end-to-end design and experimental validation of a comprehensive machine-learning framework for the analysis of network attack data that returns effective and comprehensible defence strategies for use against eithermanaged or emerging attacks. The inadequacy of classical signature-based intrusion detection methods has been acknowledged due to the serendipitous emergence of AI-powered self-learning attack methods, requiring additional investigation in the academic literature to provide preprocessing, systematic multi-algorithm comparisons, operational interpretability, and deployable system architectures.
The research involved five phases. The first phase screened through literature to identify candidate benchmark datasets, machine learning approaches, anomaly detection regimes, deep learning architectures, and hybrid ensemble systems for network intrusion detection. It was confirmed that gradient-boosting ensemble methods give the best performance on high-dimensionality tabular network traffic data, and the following four specific research gaps motivate the present work.
In phase two, the CICIDS2017 dataset was chosen as the primary experimental data source, being based on realistic generation of traffic, broad coverage of attack types, rich extraction of attack and normal features, and compelling previous work based on CICIDS2017 that allows direct comparisons to antecedent works. The dataset was acquired, combined from multiple CSV sources, and passed through a comprehensive quality assessment that identified 23{,}415 records occurring more than once, 8{,}941 records with missing or infinite values, and labels with internal inconsistencies—all of which were remediated.
In phase three, a multi-stage preprocessing pipeline was designed and executed. Z-score normalisation, log-transformations of skewed features, and casting to int 32 recommended adjustments for numerical stability of the sampled data and general-purpose normalisation of remaining numerical features. To handle the class imbalance (80.4% benign, 19.6% attack), both SMOTE augmentation and random undersampling techniques were utilised to obtain a training distribution with good recall on all attack categories. Using a three-stage feature selection procedure - Pearson correlation filtering, VIF-based multicollinearity filtering and importance pre-filtering - results in a ‘slimming’ of the features set from 80 to 74, reducing redundancy while still retaining discriminative power. The dataset from this processing is stratified into training (60%), validation (20%) and appropriately, a test (20% distribution).
The fourth phase consisted of designing, implementing and training six machine learning models: Logistic Regression, K-Nearest Neighbours, Random Forest, XGBoost, an MLP Classifier and LightGBM with each models architecture and maths documented fully. Hyperparameter optimisation is performed using 5-fold stratified cross-validation on the training set. Final performance on the held-out test set is presented using macro-averaged Accuracy / Precision / Recall / F1-score, and ROC-AUC. Finally, two additional models are employed as ‘cogs’ within a hybrid detection architecture, an unsupervised Isolation Forest model and lastly, an LSTM temporal model that is also evaluated.
In the fifth, final phase we run a global feature importance analysis of features within the model yielding the best performance, LightGBM. Feature importance with SHAP are run automatically and ‘interpret classification’ instance-level explaining with LIME. Documenting the four most discriminative features and their cybersecurity interpretation. Systematic exploration of false positives and negatives uncovered specific failure modes---notably the tendency to conflate legitimate high-volume file transfers for bandwidth exhaustion and DoS-class attacks, and the difficulty of distinguishing application-layer attacks from benign HTTP/S traffic at the flow-feature level. Most importantly, this examination generated prescriptive guidance for future improvements to the system such as the addition of rolling-variance and protocol-context features. A SIEM-targeted architecture for deploying such a system was designed, specifying where in the architecture the data is ingested, where classification occurs, what interpretability pathways are provided, and what type of integration with the SIEM is needed. The resulting scale of each component targets at network of roughly $5 \times 10^5$ flows per second on commodity server hardware.
\subsection{Principal Findings}
The experimental results support the following principal findings:
\begin{enumerate}
\item \textbf{LightGBM yields state-of-the-art performance on the CICIDS2017 benchmark:} Accuracy~=~0.98, F1-score~=~0.97, ROC-AUC~=~0.99, outperforming each of five competing algorithms while requiring the least amount of time for the entire gradient-boosting fitting procedure (89~seconds, vs 317~seconds for XGBoost). This supports the substantive value of leaf-wise tree growth, GOSS sampling, and EFB feature bundling for the network traffic classification problem.
\item \textbf{Preprocessing is key.} The shift taken from the original imbalanced dataset to the preprocessed training distribution led to qualitative improvements on minority-class recall for all models. Combining SMOTE with an undersampling approach led to a better balanced distribution than both individual techniques, and the stacked approach reduces false-positive rates of 23\u00bf for Botnet, Brute Force, and Web attack categories on average relative to models fitted on the imbalanced original data. \begin{enumerate}[label={\bfseries \arabic*.}]
\item \textbf{Feature selection improves generalisability.} By avoiding performing on the worst models trained all features, we conclude that the omitted features contributed more noise than signal. The performance improvement for the models was observed to be greatest for the KNN classifiers (an F1 improvement of 0.04) and Logistic Regression (F1 improvement of 0.03), a result consistent with these being more sensitive to irrelevant features.
\item \textbf{Interpretability exposes domain-consistent feature importance.} SHAP analysis found Flow Duration, InPktRate, OutByteRate, Bwd Packet Length Max and Source Port each to be the five most discriminative features were known to indicate malicious behaviour in the network, and importantly were interpreted by Cyber Security literature as such, imbibing our model with the confidence to decide that it is making valid comparison when activated in the current domain, i.e. will thus behave acceptably when activated live.
\item \textbf{Hybrid architectures extend zero-day detection capability.} The unsupervised mode of the Isolation Forest model achieved meaningful detection on 73% of attack instances, demonstrating a degree of zero-day detection capability worth supplementing into our supervised classifier. Additionally, the LSTM temporal model was able to improve its detection of sequential attacks (such as DDoS, PortScan) 2 percentage points over that attained by per-flow classification, illuminating the significance of temporal context in multi-step attack detection.
\item \textbf{The central research hypothesis is correct.} ML methods and data science principles applied systematically to classifying network traffic produces a system that exceeds the thresholds that were predicted for the performance measure ($F_1 > 0.95$, $\mathrm{AUC} > 0.98$) and transparent to an analyst in operations for the purpose of explainable decision making.
\end{enumerate}
\section{Scientific Contributions}
To the scientific literature, we contribute:
\begin{enumerate}[label={\bfseries \arabic*.}]
\item \textbf{A reproducible, modular preprocessing pipeline} for network intrusion detection datasets with in mind class imbalance and data quality, providing a complete, end-to-end data engineering approach which can be applied to future IDS datasets.
\item \textbf{The most systematic published comparison of six ML algorithms} on CICIDS2017 using macro-averaged metrics across all 14 traffic classes and provide per-class F1-score, Precision and Recall breakdowns.
\item \textbf{A complete SHAP and LIME interpretability analysis} that illuminating the operational applicability of explanation techniques for machine learning models amongst a particular IDS analyst workflow contexts and suggests recommendations for guard alert triage and feature engineering.
\item \textbf{Concrete SIEM-compatible deployment architecture.} We underpin our recommendations for how a future IDS might look by tackling the gap between developing our model and deploying with industry tools, dedicating our discussion to data ingestion, classification engine, interpretability module, and platform integration.
\item \textbf{Empirical evidence of Lightgbm best fitted algorithm} for network traffic multiclass classification and then why Lightgbm would be the best by virtue of the class imbalance severe, the high dimensionality of the features set uptream, and because it would cater for the requirement of an on-line architecture (no batch processing). This supplies proof to the practitioner.